\scanpage{109}

La formule ci-dessus pour

\[
\deleted{\unsure{\text{élémt}}}\quad\alpha(n(x)=\unsure{o|x}\bar x
\]

donne, si on fait $x=y+z$

\[
\begin{aligned}
\alpha(\underset{\substack{\Vert\\ n(y)+n(z)+\varphi_n(y,z)}}{\underbrace{n(y+z)}})
&=\deleted{\unsure{\alpha(n(x))}}(y+z)\overline{(y+z)}=\\
&=\underset{\substack{\Vert\\\alpha(n(y))}}{y\bar y}
+\underset{\alpha(n(z))}{z\bar z}
+\underbrace{y\bar z+\underbrace{z\bar y}_{=\overline{y\bar z}}}_{\tau(y\bar z)}
\end{aligned}
\]

\[
\text{\unsure{d'où|donc}}\qquad\alpha(\varphi_n(y,z))=\deleted{\unsure{\alpha t}}\tau(y\bar z)
\qquad\text{i.e.}\quad\alpha(t(y\bar z))
\]

donc si $\alpha$ est injectif, on en déduit

\[
\boxed{\varphi_n(y,z)=t(y\bar z)}
\]

i.e. la formule d'additivité pour $n$

\[
\boxed{n(y+z)=n(y)+n(z)+\underset{\substack{\Vert\\t(\bar yz)}}{t(y\bar z)}}
\qquad\text{car }t(x)=t(\bar x)
\]

Ceci posé, on veut

3°) la formule précédente est vraie (sans supposer $\alpha$ injective \ldots)

4°) Une formule facultative, quand $k$ est fini loc. libre \deleted{de} rang \deleted{\unsure{2}} $n$ sur un anneau $k_0$, et \deleted{alors on \unsure{ne} suppose} $A$ loc. libre de rang $2n$ sur $k_0$ (N.B. Donc $k_0\longrightarrow A$ injectif, mais on ne suppose pas $k\longrightarrow A$ injectif).

\[
\boxed{\left\{\begin{aligned}
\mathrm{Tr}_{k/k_0}(t(x))&=\mathrm{Tr}_{A/k_0}(x)\\
N_{k/k_0}(\deleted{\unsure{t}}n(x)&=N_{A/k_0}(x)
\end{aligned}\right.}
\]

\margin{identité entre applic. \uline{linéaires} $A\longrightarrow k_0$}

\margin{identité \unsure{de|des} $(2n)$-formes}

% archive sheet number: 108

